\begin{tabular}{lllcrrr}
\toprule
Class & Family & Species & IUCN & Taï NP (\%) & Mont Péko NP (\%) & Difference [95\% CI] \\
\midrule
Amphibia & Dicroglossidae & \textit{Hoplobatrachus occipitalis} & LC & 5.4 & 4.7 & +0.7 [-2.9, 4.2] \\
Aves & Bucerotidae & \textit{Lophoceros nasutus} & LC & 26.4 & 17.4 & +9.0 [2.6, 15.3] \\
 &  & \textit{Ceratogymna elata} & VU & 9.9 & 0.0 & +9.9 [6.9, 13.6] \\
 & Musophagidae & \textit{Corythaeola cristata} & LC & 21.0 & 15.2 & +5.8 [-0.1, 11.7] \\
 & Numididae & \textit{Numida meleagris} & LC & 3.6 & 0.9 & +2.7 [0.3, 5.3] \\
Gastropoda & Achatinidae & \textit{Achatina achatina} & LC & 4.8 & 0.0 & +4.8 [2.6, 7.7] \\
Mammalia & Anomaluridae & \textit{Anomalurus pelii} & LC & 6.9 & 2.2 & +4.7 [1.5, 8.1] \\
 & Bovidae & \textit{Philantomba maxwellii} & LC & 88.0 & 88.6 & -0.6 [-5.6, 4.4] \\
 &  & \textit{Tragelaphus scriptus} & LC & 80.2 & 89.2 & -9.1 [-14.5, -3.5] \\
 &  & \textit{Cephalophus niger} & LC & 60.4 & 57.3 & +3.1 [-4.5, 10.6] \\
 &  & \textit{Syncerus caffer nanus} & LC & 44.4 & 47.8 & -3.3 [-10.9, 4.3] \\
 &  & \textit{Cephalophus dorsalis} & LC & 27.0 & 32.3 & -5.3 [-12.2, 1.8] \\
 &  & \textit{Cephalophus silvicultor} & NT & 25.8 & 28.2 & -2.3 [-9.2, 4.5] \\
 &  & \textit{Tragelaphus eurycerus} & NT & 2.4 & 1.8 & +0.6 [-2.0, 3.0] \\
 & Cercopithecidae & \textit{Cercopithecus petaurista} & LC & 69.4 & 80.4 & -11.0 [-17.5, -4.3] \\
 &  & \textit{Cercopithecus campbelli} & NT & 80.0 & 65.8 & +14.2 [7.2, 20.7] \\
 &  & \textit{Cercopithecus diana} & EN & 65.2 & 62.0 & +3.1 [-4.2, 10.5] \\
 &  & \textit{Piliocolobus badius} & EN & 58.5 & 66.5 & -7.9 [-15.2, -0.4] \\
 &  & \textit{Colobus polykomos} & EN & 51.4 & 40.2 & +11.2 [3.5, 18.6] \\
 &  & \textit{Procolobus verus} & VU & 20.7 & 36.4 & -15.7 [-22.4, -8.7] \\
 & Elephantidae & \textit{Loxodonta cyclotis} & CR & 12.3 & 29.4 & -17.1 [-23.2, -10.9] \\
 & Felidae & \textit{Panthera pardus} & VU & 13.2 & 24.1 & -10.8 [-16.8, -4.9] \\
 &  & \textit{Caracal aurata} & VU & 1.8 & 1.6 & +0.2 [-2.1, 2.5] \\
 &  & \textit{Felis lybica} & LC & 2.1 & 0.9 & +1.2 [-1.0, 3.4] \\
 & Herpestidae & \textit{Galerella sanguinea} & LC & 57.4 & 73.4 & -16.1 [-23.1, -8.8] \\
 & Hippopotamidae & \textit{Choeropsis liberiensis} & EN & 3.9 & 0.0 & +3.9 [1.9, 6.6] \\
 & Hominidae & \textit{Pan troglodytes verus} & CR & 43.8 & 55.7 & -11.8 [-19.4, -4.2] \\
 & Hystricidae & \textit{Atherurus africanus} & LC & 67.9 & 88.0 & -20.1 [-26.2, -13.8] \\
 &  & \textit{Hystrix cristata} & LC & 30.9 & 53.5 & -22.5 [-29.7, -15.0] \\
 & Lorisidae & \textit{Perodicticus potto} & NT & 8.7 & 6.3 & +2.4 [-1.8, 6.5] \\
 & Manidae & \textit{Phataginus tricuspis} & EN & 65.8 & 75.0 & -9.2 [-16.1, -2.2] \\
 &  & \textit{Smutsia gigantea} & EN & 2.4 & 0.0 & +2.4 [0.7, 4.7] \\
 & Nesomyidae & \textit{Cricetomys gambianus} & LC & 55.9 & 93.3 & -37.5 [-43.3, -31.3] \\
 & Procaviidae & \textit{Dendrohyrax arboreus} & LC & 16.2 & 20.2 & -4.0 [-10.0, 1.9] \\
 &  & \textit{Procavia capensis} & LC & 0.0 & 11.1 & -11.1 [-15.0, -7.9] \\
 & Pteropodidae & \textit{Hypsignathus monstrosus} & LC & 72.1 & 73.4 & -1.3 [-8.2, 5.5] \\
 & Sciuridae & \textit{Epixerus ebii} & LC & 82.9 & 90.2 & -7.3 [-12.5, -2.0] \\
 &  & \textit{Xerus erythropus} & LC & 4.8 & 5.7 & -0.9 [-4.5, 2.6] \\
 & Suidae & \textit{Potamochoerus porcus} & LC & 29.1 & 39.2 & -10.1 [-17.3, -2.8] \\
 &  & \textit{Hylochoerus meinertzhageni} & LC & 7.5 & 11.4 & -3.9 [-8.5, 0.6] \\
 & Thryonomyidae & \textit{Thryonomys swinderianus} & LC & 96.4 & 97.2 & -0.8 [-3.7, 2.2] \\
 & Viverridae & \textit{Civettictis civetta} & NT & 44.4 & 65.5 & -21.1 [-28.3, -13.4] \\
 &  & \textit{Genetta pardina} & LC & 15.6 & 29.7 & -14.1 [-20.5, -7.7] \\
Reptilia & Crocodylidae & \textit{Osteolaemus tetraspis} & LC & 3.0 & 0.0 & +3.0 [1.2, 5.4] \\
 & Pythonidae & \textit{Python sebae} & NT & 6.3 & 0.9 & +5.4 [2.6, 8.6] \\
 & Testudinidae & \textit{Kinixys homeana} & CR & 8.1 & 1.6 & +6.5 [3.3, 10.1] \\
 & Varanidae & \textit{Varanus niloticus} & LC & 11.7 & 1.6 & +10.1 [6.5, 14.1] \\
\bottomrule
\end{tabular}
